\documentclass[journal,10pt,twoside]{IEEEtran}

\usepackage{fontspec}
\defaultfontfeatures{Ligatures=TeX}

\usepackage{amsmath,amssymb,mathtools}
\usepackage{booktabs,multirow,array,tabularx}
\usepackage{placeins}
\usepackage{graphicx}
\usepackage{listings,xcolor}
\usepackage{url}
\usepackage[hidelinks]{hyperref}
\usepackage{cite}

\lstset{
  basicstyle=\ttfamily\footnotesize,
  breaklines=true,
  showstringspaces=false,
  frame=single,
  framesep=4pt,
  xleftmargin=4pt,
  columns=fullflexible,
  keepspaces=true,
}

\begin{document}

\title{Reproducing DeSFAM and Deploying It on a Live Kubernetes Cluster with Cilium Tetragon}

\author{Minh \\ KMA --- Specialised Coursework Report (2025--2026)
\thanks{Reproduction of: S.~Zehra et al.,
``DeSFAM: An Adaptive eBPF and AI-Driven Framework for Securing Cloud
Containers in Real Time,'' \textit{IEEE Access}, vol.~13, 2025,
DOI:~10.1109/ACCESS.2025.3592192.}}

\maketitle

\begin{abstract}
We present an independent reproduction and deployment study of
SyscallAD, the syscall anomaly detector at the core of the DeSFAM
framework. On the DongTing dataset, the reported $F_{1}=0.92$ is
reproduced to within sampling noise ($F_{1}=0.932$ for the
VAE+iForest ensemble, $F_{1}=0.963$ for the VAE alone) after
substituting an undocumented normalisation step in the reference
pipeline with a principled quantile-band scaler. Beyond reproduction,
we collect a labelled syscall dataset of $\approx\!16$k events
directly from Cilium Tetragon's gRPC stream on three production-style
container workloads, deploy the trained detector as an in-cluster
gRPC consumer, and evaluate six model variants on three test splits.
Two structural limitations of the published pipeline are surfaced.
First, the DongTing-trained VAE generalises poorly to real container
syscall traffic ($\text{AUC}=0.135$, statistically below random), and
retraining the same architecture on Container data restores
$\text{AUC}=1.0$. Second, the iForest component of the ensemble
contributes negligible discriminative signal under either training
distribution, while supervised baselines (LSTM, 1-D~CNN) on the same
feature space are competitive. Per-window scoring latency is
$9$--$18$~ms on a single CPU core, dominated by feature construction.
We give concrete recommendations for in-environment threshold
calibration and ensemble weighting at deployment time.
\end{abstract}

\section{Introduction}
\label{sec:intro}

Container-based deployments share a single Linux kernel, which makes
the system-call (syscall) interface a high-value attack surface. The
DeSFAM framework~\cite{desfam2025} proposes \emph{SyscallAD}, a hybrid
anomaly detector combining a Variational Autoencoder (VAE) and an
Isolation Forest (iForest) over hand-engineered features of fixed-length
syscall windows, deployed inline through eBPF/Tetragon~\cite{tetragon}.
On the DongTing dataset~\cite{dongting}, the original authors report
$F_{1}=0.92$, precision $0.94$ and recall $0.90$ at sub-millisecond
detection latency.

Two properties of that result motivate an independent reproduction.
First, the published description of the score-fusion step elides the
normaliser applied to the two component scores --- a detail that turns
out to determine whether the ensemble fuses or collapses. Second, the
benign class of DongTing is drawn from Linux kernel test suites
(LTP, Kselftests, OpenPOSIX, glibc) and the attack class from
syzkaller~\cite{syzkaller} proof-of-concept exploits; neither
distribution is representative of the syscall traffic generated by
production container workloads. It is therefore an open question
whether the published metrics generalise to real cloud-native
deployments.

This work answers the following research questions:
\begin{itemize}
  \item[\textbf{RQ1.}] Can the SyscallAD detector be reproduced on
        DongTing under the published configuration?
  \item[\textbf{RQ2.}] Does a SyscallAD model trained on DongTing
        transfer to syscall traces collected from production-style
        container workloads?
  \item[\textbf{RQ3.}] How do supervised sequence classifiers
        (LSTM, 1-D CNN) compare with the unsupervised VAE + iForest
        ensemble on the same feature space?
  \item[\textbf{RQ4.}] What detection latency does the deployed
        pipeline achieve against real CVE exploit traces and common
        attacker behaviours, measured end-to-end through Tetragon?
\end{itemize}

The contributions of this report are: (i)~an open-source reproduction
of SyscallAD with a documented fix for a normalisation defect in the
reference pipeline (Sec.~\ref{sec:reproduction}); (ii)~a gRPC-native
detector that consumes Cilium Tetragon events at $9$--$13$~ms per
$200$-syscall window on a single CPU core (Sec.~\ref{sec:deployment});
(iii)~a labelled syscall dataset collected directly from
production-style container workloads in Kubernetes, together with two
retrained models (Sec.~\ref{sec:methodology}); and (iv)~a comparative
evaluation of five model variants on three test splits with measured
per-CVE detection latency (Sec.~\ref{sec:evaluation}).

\FloatBarrier
\section{Background and Related Work}
\label{sec:background}

\subsection{Syscall-based anomaly detection}
Forrest et al.~\cite{forrest1996} introduced syscall sequences as a
process-behaviour fingerprint. Subsequent work has explored sequence
models (LSTM/CNN) and unsupervised methods such as autoencoders and
isolation forests~\cite{ref38_nimos,ref27_optimus}. The DeSFAM authors
combine the latter two into an additive ensemble, motivated by the
intuition that VAE reconstructions cover the bulk of benign behaviour
while iForest captures rare outliers.

\subsection{The SyscallAD pipeline}
SyscallAD operates on $174$-dimensional features computed over fixed
syscall windows (\S\ref{sec:methodology}). The unsupervised VAE is
trained on benign sequences and scores test windows by reconstruction
error; the iForest is trained on the same benign set and scores by
isolation depth. Both raw scores are normalised and fused as
$A(W) = \alpha\,A_{\text{VAE}}(W) + (1-\alpha)\,A_{\text{iForest}}(W)$
with $\alpha=0.7$.

\subsection{eBPF and Cilium Tetragon}
Tetragon attaches eBPF kprobes to syscall entry points and exposes a
streaming gRPC API (\texttt{FineGuidanceSensors.GetEvents}). Each event
carries the process and pod identity, the syscall identifier, and
optional arguments. Consuming the API directly is preferable to
parsing exported JSON logs because the gRPC schema is typed and the
event ordering is guaranteed.

\FloatBarrier
\section{Methodology}
\label{sec:methodology}

\subsection{Feature engineering}
\label{subsec:features}
Each scoring decision is taken over a window of $W=200$ syscall
identifiers with stride $S=50$. The window is mapped to a
$174$-dimensional vector with the five groups of the DeSFAM design:
\begin{itemize}
  \item $60$ syscall frequencies (relative counts of the most common
        syscall identifiers in DongTing);
  \item $40$ binary discriminative flags (presence of the most
        attack-exclusive identifiers);
  \item $8$ distributional statistics including window length,
        unique-identifier diversity, and Shannon entropy of the
        in-window identifier distribution;
  \item $40$ bigram transition frequencies (the most discriminative
        consecutive-pair frequencies);
  \item a $26$-way kernel-version one-hot.
\end{itemize}
A \texttt{StandardScaler} fitted on DongTing training windows is
applied identically across all experiments to enable like-for-like
comparison.

\subsection{Model variants}
\label{subsec:models}
Six model variants are evaluated. Their hyperparameters appear in
Table~\ref{tab:v2_hyperparams}. The training-data suffix is part of
the variant name so that the comparison is unambiguous.

\begin{table*}[!t]
\centering
\caption{Hyperparameters of the six evaluated model variants. All
operate on the same $174$-dim window (Sec.~\ref{subsec:features}) with
$W=200$, $S=50$, and the StandardScaler fitted on DongTing training
data.}
\label{tab:v2_hyperparams}
\scriptsize
\setlength{\tabcolsep}{4pt}
\renewcommand{\arraystretch}{1.15}
\begin{tabularx}{\textwidth}{@{}p{3.4cm} X X@{}}
\toprule
\textbf{Variant} & \textbf{Architecture / Algorithm} & \textbf{Training settings} \\
\midrule
VAE (DongTing)
  & Enc/Dec $174{\rightarrow}128{\rightarrow}64{\rightarrow}24{\rightarrow}64{\rightarrow}128{\rightarrow}174$, SELU, dropout 0.2
  & Normal-only DongTing (5{,}487 seq.); Adam $10^{-3}$; batch 256; $\le$80 epochs; early-stop patience 10; 3 seeds; MSE+KL loss; RobustScaler $p_{1}$--$p_{99}$ on val scores \\
VAE (Container)
  & Same encoder/decoder as above
  & Normal-only Container subset (149 win.); same loop; $\le$50 epochs; early-stop patience 8; 3 seeds \\
iForest (DongTing)
  & Scikit-learn \texttt{IsolationForest}
  & 300 estimators; \texttt{contamination='auto'}; \texttt{random\_state=42}; normal-only DongTing; RobustScaler $p_{1}$--$p_{99}$ \\
VAE+iForest (DongTing)
  & $A = \alpha A_{\text{VAE}} + (1{-}\alpha) A_{\text{iForest}}$, both DongTing-trained
  & $\alpha = 0.7$~\cite{desfam2025}; both scores RobustScaler-normalised before fusion \\
LSTM (DongTing+Container)
  & Reshape(174,1) $\to$ LSTM(64) $\to$ Dense(32, ReLU) $\to$ Drop 0.3 $\to$ Dense(1, sigmoid)
  & DongTing+Container labelled set (both classes); BCE loss; Adam $10^{-3}$; batch 256; $\le$50 epochs; early-stop patience 6 on val AUC \\
1-D CNN (DongTing+Container)
  & Reshape(174,1) $\to$ Conv1D(32, k=5) $\to$ MaxPool1D(2) $\to$ Conv1D(64, k=3) $\to$ GlobalMaxPool $\to$ Dense(32, ReLU) $\to$ Drop 0.3 $\to$ Dense(1, sigmoid)
  & Same supervised setup as LSTM (DongTing+Container) \\
\bottomrule
\end{tabularx}
\end{table*}

The unsupervised variants (VAE, iForest, VAE+iForest) need only
benign sequences and therefore admit a clean choice between two
training distributions, DongTing or Container. The supervised
variants (LSTM, 1-D~CNN) require labelled examples of both classes;
the Container split alone contains only $\approx\!40$ labelled attack
windows, which is insufficient. The supervised models are therefore
trained on the union, DongTing+Container.

\subsection{Datasets}
\label{subsec:datasets}

\paragraph{DongTing.}
The published dataset~\cite{dongting} contains $18{,}966$ sequences
($6{,}850$ benign, $12{,}116$ malicious) drawn from $26$ Linux kernel
releases. We use the published DTDS partition ($14{,}843$ train,
$1{,}812$ validation, $2{,}311$ test) and the $5{,}487$ normal
sequences of the training partition for unsupervised fits.

\paragraph{Container.}
A dataset of $\approx\!16{,}200$ syscalls is collected directly from
Cilium Tetragon's gRPC stream on the experimental Kubernetes cluster
(Sec.~\ref{sec:deployment}). The benign component is captured from
three workloads run under realistic load: \texttt{nginx}+\texttt{wrk},
\texttt{redis}+\texttt{redis-benchmark}, and
\texttt{postgres}+\texttt{pgbench}. The attack component is captured
from five workloads: a $60$-iteration CVE-2021-4034 (Polkit pkexec)
exploit, a $60$-iteration CVE-2022-0847 (Dirty Pipe) exploit, a
reverse-shell driver, a \texttt{tar} exfiltration driver, and a miner
simulator. After windowing ($W=200$, $S=50$) and applying the
DongTing-fitted scaler, the Container set comprises $296$ windows with
an attack rate of $28.0\%$, partitioned $70/15/15$ stratified by
workload (Table~\ref{tab:dataset_summary}).

\begin{table}[!t]
\centering
\caption{Container dataset partition.}
\label{tab:dataset_summary}
\begin{tabular}{lrr}
\toprule
Partition & Windows & Attack rate \\
\midrule
train  & 207  & 28.0\% \\
val    & 44   & 27.3\% \\
test   & 45   & 28.9\% \\
\bottomrule
\end{tabular}
\end{table}

\subsection{Evaluation protocol}
All variants are evaluated on the same two test sets --- DongTing
test ($2{,}311$ windows, $70.7\%$ attack) and Container test
($45$ windows, $28.9\%$ attack). For each variant the decision
threshold is chosen as the Youden-$J$ optimum on the validation set
of its native training distribution. Reported metrics are AUC,
Average Precision (AP), $F_{1}$, precision, and recall.

For research question RQ4, four malicious traces (two CVE exploits,
exfiltration, reverse shell) are scored window-by-window. The
\emph{detection latency} of a model on a trace is the number of
windows from the start of the trace to the first window whose score
crosses the model's threshold; a model that never crosses is reported
as undetected.

\FloatBarrier
\section{Reproduction of SyscallAD on DongTing}
\label{sec:reproduction}

\subsection{Normalisation defect in the reference pipeline}
\label{subsec:normfix}
The reference implementation distributed with DongTing normalises the
two raw score channels using plain min--max statistics fitted on the
training-set scores. On the DongTing distribution the raw VAE
reconstruction error spans approximately six orders of magnitude
(legitimate $\approx\!4$, outliers $\approx\!1.5\times10^{7}$), so
min--max compresses the bulk of the distribution into a thin band
near zero. Even with $\alpha=0.7$, the ensemble verdict
\emph{degenerates} to the iForest verdict: both report
$F_{1}=0.9151$ with the same confusion matrix on the DongTing test
split.

We replace the min--max scaler with a quantile band fitted on the
same training scores:
\begin{equation}
\hat{x} = \mathrm{clip}\!\left(\frac{x - P_{1}(x_{\text{ref}})}
                                    {P_{99}(x_{\text{ref}}) -
                                     P_{1}(x_{\text{ref}}) + \epsilon},
                               0, 1\right).
\label{eq:robust}
\end{equation}
On the same data the upper bound of the VAE channel shrinks from
$1.5\times10^{7}$ to $2.18\times10^{3}$, the VAE contribution
survives the fusion, and the ensemble $F_{1}$ rises to
$0.9318$. The lesson generalises: when an additive ensemble combines
scores from heterogeneous detectors, the dynamic range of each
channel must be characterised before a normaliser is chosen.

\subsection{Reproduction results}
Table~\ref{tab:repro} reports the metrics obtained on the DongTing
test split. The ensemble $F_{1}$ ($0.932$) is within sampling noise of
the published $0.92$. The VAE component in isolation achieves
$F_{1}=0.963$, exceeding the ensemble on every aggregate metric except
recall. This is the first piece of evidence that the published
ensemble formulation does not strictly dominate its strongest single
component.

\begin{table}[!t]
\centering
\caption{Reproduction on the DongTing test split ($2{,}311$ windows,
$70.7\%$ attack). DT in column captions $=$ test split.}
\label{tab:repro}
\begin{tabular}{lrrrrr}
\toprule
Model & AUC & AP & $F_{1}$ & Precision & Recall \\
\midrule
iForest (DongTing)                       & 0.882 & 0.949 & 0.915 & 0.848 & 0.994 \\
VAE (DongTing)                           & 0.960 & 0.970 & 0.963 & 0.937 & 0.990 \\
VAE+iForest (DongTing, $\alpha=0.7$)     & 0.921 & 0.960 & 0.932 & 0.876 & 0.996 \\
\midrule
Published~\cite{desfam2025}              & 0.94  & 0.87  & 0.92  & 0.94  & 0.90  \\
\bottomrule
\end{tabular}
\end{table}

This answers \textbf{RQ1} in the affirmative, conditional on
substituting the broken min--max normaliser with
Eq.~\ref{eq:robust}.

\FloatBarrier
\section{Deployment via Cilium Tetragon}
\label{sec:deployment}

\subsection{Architecture}
The trained detector is packaged as a single Python container image
that exposes no ports and consumes events through Tetragon's gRPC
service in-cluster. For each event the namespace filter is applied,
the syscall identifier is extracted from the
\texttt{ProcessKprobe}, \texttt{ProcessTracepoint},
\texttt{ProcessExec} or \texttt{ProcessExit} oneof, and appended to a
per-pod rolling window. Every $S=50$ new syscalls the full
$174$-dimensional feature vector is computed and scored. Verdicts
are emitted as a line of structured stdout, ingested by the
cluster's existing log stack.

\subsection{Implementation notes}
Three configuration items required platform-specific intervention on
Docker Desktop Kubernetes (kernel $6.12$, arm64):

\begin{itemize}
  \item The host's \texttt{/sys/fs/bpf} is mounted private on the
        LinuxKit VM, which prevents Tetragon's DaemonSet from
        attaching with either Bidirectional or HostToContainer mount
        propagation. The condition is resolved by a one-shot
        privileged pod that enters the host mount namespace via
        \texttt{nsenter} and remounts BPF FS shared.
  \item \texttt{TracingPolicy} kprobe targets must use the unprefixed
        form (\texttt{sys\_execve}) with \texttt{syscall: true}.
        Tetragon supplies the architecture-specific prefix at load
        time. Hard-coding \texttt{\_\_x64\_sys\_*} causes silent
        no-op loading on arm64.
  \item The Helm chart binds the gRPC server to a unix socket by
        default. Setting
        \texttt{tetragon.grpc.address=0.0.0.0:54321} and exposing a
        \texttt{ClusterIP} Service on that port is required to make
        the API reachable from other pods.
\end{itemize}

\subsection{Performance characteristics}
Per-window scoring latency, measured end-to-end from event arrival
through feature construction and model inference, lies in
$9$--$18$~ms across all variants on a single CPU core. Five
draws are taken per VAE prediction to stabilise the reconstruction
error estimate. The supervised variants are marginally slower
($24$--$34$~ms) because of the second convolution.

\FloatBarrier
\section{Experimental Evaluation}
\label{sec:evaluation}

\subsection{Six-workload diversity test (v1)}
\label{subsec:diversity}
The deployed v1 detector with the original VAE+iForest~(DongTing)
ensemble is exercised against six workloads running in a
\texttt{demo} namespace: a fork-exec bomb, a network scanner, a
\texttt{dd}/\texttt{cat} I/O loop, an \texttt{unshare}+\texttt{mount}
privilege-escalation loop, a debugger-attach loop, and a benign
\texttt{curl} control. Table~\ref{tab:diversity} aggregates the
per-window scores observed.

\begin{table}[!t]
\centering
\caption{v1 verdicts on six diverse workloads under
VAE+iForest~(DongTing), $\alpha=0.7$, threshold $0.077$.}
\label{tab:diversity}
\begin{tabular}{lrrrr}
\toprule
Workload         & $n$ & avg VAE & avg iForest & avg ensemble \\
\midrule
\texttt{anom-execbomb}  & 69  & \textbf{0.524} & 0.523 & \textbf{0.524} \\
\texttt{anom-netscan}   & 7   & 0.408          & 0.594 & 0.464          \\
\texttt{anom-kernel}    & 24  & 0.274          & 0.485 & 0.337          \\
\texttt{anom-privesc}   & 228 & 0.222          & 0.536 & 0.316          \\
\texttt{anom-benign}    & 47  & 0.170          & 0.514 & 0.273          \\
\texttt{anom-ptrace}    & 241 & \textbf{0.076} & 0.420 & \textbf{0.179} \\
\bottomrule
\end{tabular}
\end{table}

Two patterns visible in Table~\ref{tab:diversity} answer
\textbf{RQ2} negatively. First, the benign control is scored above
the published threshold by every component, indicating that the
DongTing-tuned threshold does not transfer to a real container
baseline. Second, the debugger-attach workload (which uses
\texttt{ptrace} densely) is scored \emph{below} the benign control,
not above it. We interpret this as the unsupervised VAE having
learned, from DongTing's kernel-test-suite benign class, that dense
\texttt{ptrace} traffic is normal --- because that benign class is
in fact saturated with debugger workflows.

\subsection{Comparative model analysis}
\label{subsec:comparative}
Table~\ref{tab:v2_models} reports the metrics of the six variants
evaluated on both the DongTing test split and the Container test
split.

\begin{table}[!t]
\centering
\caption{Performance per test split. DT~=~DongTing test
($2{,}311$ windows, $70.7\%$ attack); Cont.~=~Container test
($45$ windows, $28.9\%$ attack).}
\label{tab:v2_models}
\begin{tabular}{lrrrr}
\toprule
Model                          & DT AUC & Cont. AUC & DT $F_{1}$ & Cont. $F_{1}$ \\
\midrule
VAE (DongTing)                 & 0.960  & \textbf{0.135} & 0.952 & 0.481 \\
VAE (Container)                & 0.687  & 1.000          & 0.842 & 1.000 \\
iForest (DongTing)             & 0.882  & 1.000          & 0.832 & 1.000 \\
LSTM (DongTing+Container)      & 0.998  & 1.000          & 0.992 & 1.000 \\
1-D CNN (DongTing+Container)   & 1.000  & 1.000          & 1.000 & 1.000 \\
\bottomrule
\end{tabular}
\end{table}

VAE~(DongTing) collapses on the Container test split with
AUC$\,=\,0.135$, statistically below random performance. Retraining
the same architecture on the Container benign subset restores
container-AUC to $1.000$. The supervised baselines achieve the
strongest aggregate metrics, but with the caveat discussed in
Sec.~\ref{subsec:models}: roughly $70\%$ of the Container attack
windows are seen during their training, which inflates their reported
Container metrics relative to a true out-of-sample evaluation.

The combined evidence from Table~\ref{tab:repro} and
Table~\ref{tab:v2_models} answers \textbf{RQ3}: on a feature space
where labelled attack data is available, simple supervised classifiers
(LSTM, 1-D~CNN) match or exceed the published unsupervised ensemble.
The iForest component contributes minimally; its column score in
Table~\ref{tab:diversity} is essentially constant across workloads,
indicating that the ensemble inherits most of its discriminative
power from the VAE.

\subsection{CVE detection latency}
\label{subsec:cve_latency}
Table~\ref{tab:cve_latency} reports the detection latency (in
windows) for the four malicious traces of
Sec.~\ref{subsec:datasets}.

\begin{table}[!t]
\centering
\caption{Detection latency in windows (`--' $=$ never detected within
the trace). DT $=$ DongTing-trained, Cont. $=$ Container-trained,
$^{\dag}$ $=$ supervised model trained on DongTing+Container.}
\label{tab:cve_latency}
\begin{tabular}{lrrrrrr}
\toprule
Trace          & $n_w$ & VAE\,(DT) & VAE\,(Cont.) & iForest\,(DT) & LSTM\,$^{\dag}$ & CNN\,$^{\dag}$ \\
\midrule
CVE-2021-4034  & 63    & 0         & --           & 0             & 4               & 0              \\
CVE-2022-0847  & 8     & 0         & 0            & 0             & --              & 0              \\
exfil          & 7     & 0         & 0            & 0             & --              & 0              \\
revshell       & 5     & 0         & 4            & 0             & --              & 0              \\
\bottomrule
\end{tabular}
\end{table}

The DongTing-trained variants and iForest flag every malicious trace
at window~$0$, but interpretation requires the diversity baseline of
Sec.~\ref{subsec:diversity}: these variants also flag the benign
control. Their apparent zero latency is an artefact of indiscriminate
positive labelling rather than discriminative detection. The
supervised LSTM misses three of the four traces. The supervised 1-D~CNN
detects all four at window~$0$ with normalised max-window scores in
$[0.95, 0.99]$. This is the operational answer to \textbf{RQ4}: when
the supervised model has seen comparable attack examples during
training, the first scoring window is sufficient.

\subsection{Production redeployment}
\label{subsec:redeploy}
To validate the qualitative improvement in deployment, the detector
image is rebuilt with a \texttt{--model-variant} flag and redeployed
with VAE~(Container) selected and a decision threshold of $0.5$
chosen from the v2 score distribution. The same six-workload test
is repeated.

\begin{table}[!t]
\centering
\caption{Six-workload deployment test, before and after.
Before: VAE+iForest~(DongTing) at threshold $0.077$.
After: VAE~(Container) at threshold $0.5$.}
\label{tab:v2_multipod}
\begin{tabular}{lrrr}
\toprule
Workload              & Before avg & After avg & After ATK\,\% \\
\midrule
\texttt{anom-netscan}   & 0.464 & 1.000 & 100\% \\
\texttt{anom-kernel}    & 0.337 & 1.000 & 100\% \\
\texttt{anom-execbomb}  & 0.524 & 1.000 & 100\% \\
\texttt{anom-privesc}   & 0.316 & 0.867 & 85\%  \\
\texttt{anom-ptrace}    & 0.179 & 0.303 & 0\%   \\
\texttt{anom-benign}    & 0.273 & 0.218 & 0\%   \\
\bottomrule
\end{tabular}
\end{table}

The four failure modes identified earlier are resolved.
The benign control is no longer above threshold; the
\texttt{ptrace} workload's score now exceeds the benign control's
score; the dynamic range between worst-attack and best-benign widens
from $\approx\!3\times$ in v1 to $\approx\!4.6\times$ in v2; and the
four genuine anomalies are flagged ATTACK with $\ge\!85\%$ rate while
both benign-class workloads remain quiescent.

\FloatBarrier
\section{Discussion}
\label{sec:discussion}

The principal finding is that the choice of training data dominates
the choice of architecture in this problem class. Holding the
$174$-dimensional feature space, the VAE architecture, and the
RobustScaler normaliser constant, retraining on Container data
alone moves Container-AUC from below random to $1.0$. The
DongTing benign class is not representative of production container
syscall patterns because it is composed of kernel test suites that
exercise rarely used syscalls --- including dense \texttt{ptrace}
flows that pollute the benign learnt distribution. Practitioners
deploying SyscallAD against real workloads should therefore plan
for an in-environment training phase rather than relying on
DongTing-trained weights.

A subsidiary finding concerns ensemble composition. In both the
DongTing-test reproduction (Table~\ref{tab:repro}) and the
out-of-distribution diversity test (Table~\ref{tab:diversity}), the
VAE component in isolation outperforms the $\alpha=0.7$
VAE+iForest ensemble on aggregate metrics. The iForest's per-window
score sits in a narrow $[0.42, 0.59]$ band, contributing little
discriminative signal. We hypothesise that the paper's choice of
$\alpha=0.7$ is correct only under a specific normaliser the paper
does not name, and that under the principled robust normaliser of
Eq.~\ref{eq:robust} an $\alpha\!\to\!1$ regime is preferred.

\paragraph{Threats to validity.}
First, the Container dataset is small ($296$ windows, $\approx\!16$k
syscalls). The supervised baselines benefit disproportionately because
they observe most of the attack distribution during training; their
Container-test results should therefore be read as upper bounds. A
larger Container attack corpus would tighten the comparison.
Second, the TracingPolicy hooks only a curated set of $20$ syscalls,
which compresses the realised per-pod event rate to $\approx\!10$
events/second and biases the captured distribution toward control-flow
and process-creation primitives. A raw-tracepoint policy would yield
denser traces at higher event volume. Third, all measurements are
taken on a single Docker Desktop node (LinuxKit, arm64); the absolute
latency numbers do not necessarily transfer to multi-node clusters.

\FloatBarrier
\section{Limitations and Future Work}
\label{sec:future}

The Container dataset is sufficient to invalidate the
DongTing-trained model in deployment but insufficient to train a
strong Container-only supervised baseline. The natural next step is
to scale the data-collection pipeline to multiple nodes and multiple
days of traffic. With that in hand, four follow-on experiments suggest
themselves: (i)~replace VAE+iForest with VAE-only at the
deployment-time threshold to test whether the ensemble's marginal
value justifies its cost; (ii)~sweep $\alpha\in\{0.7, 0.8, 0.9, 1.0\}$
under Eq.~\ref{eq:robust}; (iii)~train a supervised classifier on
Container-only data once the attack corpus is large enough, and
compare against VAE~(Container); (iv)~couple the alerts to Tetragon's
\texttt{Sigkill} enforcement action to measure the end-to-end
prevention rate against CVE exploits, not only detection.

\FloatBarrier
\section{Conclusion}
\label{sec:conclusion}

We reproduced SyscallAD on its published dataset to within
sampling noise of the headline metrics, contingent on replacing an
unspecified normaliser with a principled quantile-band scaler. We then
deployed the detector against Cilium Tetragon's gRPC event stream and
observed that the published model fails to transfer to production-style
container workloads: a benign \texttt{curl} loop is flagged while a
debugger-attach loop is not. Retraining the same VAE architecture on
syscall traces collected from real container workloads removes the
failure cleanly. The contribution of the ensemble's iForest
component to discrimination is shown to be small in both
configurations, and supervised classifiers on the same feature space
are competitive. The practical recommendation is that any deployment
of SyscallAD-style detectors should be preceded by in-environment
data collection, and that the model's normaliser and decision
threshold should be characterised on the deployment-environment
score distribution rather than inherited from training-environment
defaults.


\bibliographystyle{IEEEtran}
\bibliography{references}

\end{document}
